\documentclass[aspectratio=169]{beamer}

% ── Kodowanie i język ──
\usepackage[utf8]{inputenc}
\usepackage[T1]{fontenc}
\usepackage[polish]{babel}

% ── Pakiety ──
\usepackage{graphicx}
\usepackage{listings}
\usepackage{xcolor}
\usepackage{tikz}
\usepackage{fontawesome5}
\usepackage{hyperref}

\usetikzlibrary{arrows.meta, positioning, shapes.geometric, fit, calc}

% ── Motyw ──
\usetheme{Madrid}
\usecolortheme{whale}

\definecolor{claudepurple}{RGB}{100, 65, 165}
\definecolor{codebg}{RGB}{30, 30, 46}
\definecolor{titlebarbg}{RGB}{50, 50, 66}
\definecolor{codegreen}{RGB}{80, 200, 120}
\definecolor{codeyellow}{RGB}{250, 189, 47}
\definecolor{codered}{RGB}{243, 95, 116}
\definecolor{codegray}{RGB}{130, 130, 160}
\definecolor{promptgreen}{RGB}{166, 227, 161}

% ── Makra pomocnicze ──
\newcommand{\separator}{\color{codegray}\rule{\linewidth}{0.3pt}}
\newcommand{\toolicon}{\color{codegray}\faCog\ }
\newcommand{\toolok}{\color{codegreen}\faCheck\ }
\newcommand{\toolfail}{\color{codered}\faTimes\ }
\newcommand{\infoicon}{\color{codegray}\faInfoCircle\ }
\newcommand{\bannerline}{\color{cyan}\rule{\linewidth}{0.5pt}}

% ── Makro terminala ──
\newcommand{\terminal}[3]{% {szerokość}{tytuł}{treść}
\begin{tikzpicture}
    \node[
        fill=codebg,
        rounded corners=4pt,
        inner sep=0pt,
        minimum width=#1,
        draw=titlebarbg,
        line width=0.5pt
    ] (term) {%
        \begin{minipage}{#1}
            % pasek tytułowy
            \begin{tikzpicture}
                \fill[titlebarbg, rounded corners=4pt]
                    (0,0) rectangle (\linewidth, 0.45cm);
                % nie zaokrąglaj dołu
                \fill[titlebarbg]
                    (0,0) rectangle (\linewidth, 0.15cm);
                % kropki macOS
                \fill[codered] (0.3,0.225) circle (0.08cm);
                \fill[codeyellow] (0.65,0.225) circle (0.08cm);
                \fill[codegreen] (1.0,0.225) circle (0.08cm);
                % tytuł
                \node[anchor=center, font=\tiny\sffamily, text=codegray]
                    at (0.5\linewidth, 0.225) {#2};
            \end{tikzpicture}\\[-2pt]
            % treść
            \vspace{4pt}
            \hspace{6pt}\begin{minipage}{\dimexpr#1-14pt}
                \color{white}\ttfamily\tiny\raggedright
                #3
            \end{minipage}
            \vspace{6pt}
        \end{minipage}
    };
\end{tikzpicture}
}

\setbeamercolor{title}{fg=white, bg=claudepurple}
\setbeamercolor{frametitle}{fg=white, bg=claudepurple}
\setbeamercolor{structure}{fg=claudepurple}
\setbeamercolor{block title}{bg=claudepurple, fg=white}

% ── Listings ──
\lstset{
    basicstyle=\ttfamily\scriptsize,
    backgroundcolor=\color{codebg},
    commentstyle=\color{codegreen},
    keywordstyle=\color{cyan},
    stringstyle=\color{orange},
    breaklines=true,
    frame=single,
    rulecolor=\color{claudepurple},
    showstringspaces=false
}

% ── Metadane ──
\title{Claude Code Python}
\subtitle{Terminalowy asystent AI do programowania}
\author{Zastosowanie Metod Sztucznej Inteligencji}
\date{\today}
\institute{}

% ==========================================================================
\begin{document}

% ══════════════════════════════════════════════════════════════════════════
% SLAJD 1 — Tytuł
% ══════════════════════════════════════════════════════════════════════════
\begin{frame}
    \titlepage
\end{frame}

% ══════════════════════════════════════════════════════════════════════════
% SLAJD 2 — Cel projektu + technologie
% ══════════════════════════════════════════════════════════════════════════
\begin{frame}{O~projekcie}
    \begin{columns}[T]
        \column{0.55\textwidth}
        \begin{block}{Cel}
            Terminalowy asystent AI wykorzystujący Anthropic API
            do autonomicznego programowania --- czyta, edytuje pliki
            i~wykonuje komendy systemowe.
        \end{block}

        \vspace{0.2cm}
        \textbf{Kluczowe funkcje:}
        \begin{itemize}\setlength{\itemsep}{1pt}
            \item \faRobot\ Konwersacja z~modelem Claude
            \item \faTools\ Tool calling (6~narzędzi)
            \item \faBrain\ Extended thinking (4~poziomy)
            \item \faUsersCog\ System 4~agentów
            \item \faTerminal\ Tryb REPL + CLI
        \end{itemize}

        \column{0.42\textwidth}
        \begin{block}{Technologie}
            \begin{itemize}\setlength{\itemsep}{1pt}
                \item \textbf{Python 3.8+}
                \item \textbf{anthropic} --- SDK
                \item \textbf{Textual} --- TUI
                \item \textbf{httpx} --- streaming
                \item \textbf{Flask} --- OAuth proxy
            \end{itemize}
        \end{block}

        \vspace{0.2cm}
        \begin{block}{API}
            \texttt{/v1/messages}\\
            Model: \texttt{claude-sonnet-4-5}\\
            Auth: API Key / OAuth PKCE
        \end{block}
    \end{columns}
\end{frame}

% ══════════════════════════════════════════════════════════════════════════
% SLAJD 3 — Zastosowanie: Interaktywny czat + edycja plików
% ══════════════════════════════════════════════════════════════════════════
\begin{frame}{Czat z~modelem i~edycja plików}
    \begin{columns}[T]
        \column{0.48\textwidth}
        \terminal{0.95\columnwidth}{claude-code-py}{%
            \bannerline\\
            \color{cyan}Claude Code Python - Sonnet 4.5\\
            \color{cyan}Extended Thinking Enabled\\
            \bannerline\\[3pt]
            \infoicon\color{white}Model: claude-sonnet-4-5\\
            \infoicon\color{white}Tools enabled (6 tools)\\
            \separator\\[2pt]
            \color{promptgreen}\textbf{You:} \color{white}Wyja\'snij wzorzec Observer\\
            \separator\\[2pt]
            \color{codegray}$\therefore$ Thinking...\\
            \color{codegray}$\therefore$ Thinking complete (1s, 60 tokens)\\[2pt]
            \color{cyan}\textbf{Claude:} \color{white}Wzorzec Observer to\\
            behawioralny wzorzec projektowy,\\
            w~kt\'orym obiekt utrzymuje list\k{e}\\
            obserwator\'ow i~powiadamia ich\\
            o~ka\.zdej zmianie stanu.%
        }

        \vspace{0.15cm}
        \begin{itemize}\setlength{\itemsep}{0pt}\footnotesize
            \item Pe\l{}noekranowy TUI (Textual)
            \item Streaming w~czasie rzeczywistym
            \item Widoczny proces my\'slenia
        \end{itemize}

        \column{0.48\textwidth}
        \terminal{0.95\columnwidth}{claude-code-py}{%
            \color{promptgreen}\textbf{You:} \color{white}Napraw buga w~calculate\_total()\\
            \separator\\[2pt]
            \color{codegray}$\therefore$ Thinking complete (1s, 68 tokens)\\[2pt]
            \toolicon Executing tool:\color{white}read\_file\\
            \color{codegray}\quad Input: \color{white}src/billing.py\\
            \toolok\color{white}Tool result (success)\\[2pt]
            \toolicon Executing tool:\color{white}edit\_file\\
            \color{codegray}\quad Input: \color{white}linia 42\\
            \toolok\color{white}Tool result (success)\\[2pt]
            \color{cyan}\textbf{Claude:} \color{white}Brakowa\l{}o sprawdzenia\\
            czy \texttt{discount} nie jest \texttt{None}.%
        }

        \vspace{0.15cm}
        \begin{itemize}\setlength{\itemsep}{0pt}\footnotesize
            \item Model sam czyta, szuka, edytuje
            \item 6~narz\k{e}dzi: Bash, Read, Write, Edit, Grep, Glob
            \item P\k{e}tla do 5~rund automatycznie
        \end{itemize}
    \end{columns}
\end{frame}

% ══════════════════════════════════════════════════════════════════════════
% SLAJD 4 — Zastosowanie: Czytanie i edycja plików (wierszyk)
% ══════════════════════════════════════════════════════════════════════════
\begin{frame}{Czytanie i~edycja plik\'ow}
    \begin{columns}[T]
        \column{0.48\textwidth}
        \terminal{0.95\columnwidth}{claude-code-py}{%
            \color{promptgreen}\textbf{You:} \color{white}widzisz wierszyk?\\
            zmie\'n wers 4 i~napisz na nowo\\
            drug\k{a} strof\k{e}\\
            \separator\\[2pt]
            \color{codegray}$\therefore$ Thinking complete (1s, 56)\\[1pt]
            \toolicon Executing tool:\color{white}bash\\
            \color{codegray}\quad \color{white}Get-ChildItem\\
            \toolok\color{white}Tool result (success)\\[1pt]
            \color{codegray}$\therefore$ Thinking complete (1s, 38)\\[1pt]
            \toolicon Executing tool:\color{white}read\_file\\
            \color{codegray}\quad \color{white}wierszyk.txt\\
            \toolok\color{white}Tool result (success)\\[1pt]
            \color{codegray}$\therefore$ Thinking complete (2s, 102)\\[1pt]
            \toolicon Executing tool:\color{white}write\_file\\
            \toolok\color{white}Tool result (success)%
        }

        \column{0.48\textwidth}
        \terminal{0.95\columnwidth}{claude-code-py}{%
            \color{cyan}\textbf{Claude:} \color{white}Gotowe! Zmieni\l{}em\\
            wers 4 i~napisa\l{}em now\k{a} strof\k{e}:\\[3pt]
            \color{codegray}\textit{Na biurku siedzi rudy kot,}\\
            \color{codegray}\textit{Patrzy w~ekran uwa\.znie.}\\
            \color{codegray}\textit{\L{}apk\k{a} dotyka klawiatury,}\\
            \color{white}\textit{Debuguje p\k{e}tle i~tablice.}\\[3pt]
            \color{white}\textit{W~testach szuka b\l{}\k{e}d\'ow tkwi,}\\
            \color{white}\textit{Przegl\k{a}da logi uwa\.znie.}\\
            \color{white}\textit{Refaktoryzuje kod stary,}\\
            \color{white}\textit{I~komituje zadowolony.}%
        }

        \vspace{0.15cm}
        \begin{itemize}\setlength{\itemsep}{0pt}\footnotesize
            \item Jedno polecenie $\rightarrow$ 3~narz\k{e}dzia
            \item Model sam znajduje, czyta, edytuje
            \item Kreatywna modyfikacja tre\'sci
        \end{itemize}
    \end{columns}
\end{frame}

% ══════════════════════════════════════════════════════════════════════════
% SLAJD 5 — Zastosowanie: Debugowanie + komendy systemowe
% ══════════════════════════════════════════════════════════════════════════
\begin{frame}{Autonomiczne naprawianie bug\'ow}
    \begin{columns}[T]
        \column{0.48\textwidth}
        \terminal{0.95\columnwidth}{claude-code-py}{%
            \color{promptgreen}\textbf{You:} \color{white}Napraw buga\\
            w~calculate\_total()\\
            \separator\\[2pt]
            \color{codegray}$\therefore$ Thinking complete (1s, 85)\\[1pt]
            \toolicon Executing tool:\color{white}glob\\
            \color{codegray}\quad \color{white}*.py\\
            \toolok\color{white}Tool result (success)\\[1pt]
            \color{codegray}$\therefore$ Thinking complete (1s, 29)\\[1pt]
            \toolicon Executing tool:\color{white}grep\\
            \color{codegray}\quad \color{white}calculate\_total\\
            \toolok\color{white}Tool result (success)\\[1pt]
            \toolicon Executing tool:\color{white}read\_file\\
            \color{codegray}\quad \color{white}demo\_billing.py\\
            \toolok\color{white}Tool result (success)\\[1pt]
            \toolicon Executing tool:\color{white}edit\_file\\
            \toolok\color{white}Tool result (success)\\[1pt]
            \toolicon Executing tool:\color{white}read\_file\\
            \toolok\color{white}Tool result (success)%
        }

        \column{0.48\textwidth}
        \terminal{0.95\columnwidth}{claude-code-py}{%
            \color{cyan}\textbf{Claude:} \color{white}Bug Fixed!\\[3pt]
            \color{codered}Root cause: \color{white}\texttt{None} check\\
            \color{white}\texttt{total - discount} crashes\\
            when \texttt{discount=None}\\[3pt]
            \color{codegreen}Fix:\color{white}\\
            \texttt{if discount is not None:}\\
            \texttt{~~~~total = total - discount}\\[3pt]
            \color{codegray}+ test cases\\
            + prevention strategies%
        }

        \vspace{0.15cm}
        \begin{itemize}\setlength{\itemsep}{0pt}\footnotesize
            \item Szuka $\rightarrow$ czyta $\rightarrow$ naprawia $\rightarrow$ weryfikuje
            \item Agent Bug Fixer wzmacnia prompt
            \item 5~narz\k{e}dzi w~jednym zapytaniu
        \end{itemize}
    \end{columns}
\end{frame}

% ══════════════════════════════════════════════════════════════════════════
% SLAJD 6 — Zastosowanie: Extended Thinking + Agenci
% ══════════════════════════════════════════════════════════════════════════
\begin{frame}{Extended thinking i~agenci}
    \begin{columns}[T]
        \column{0.48\textwidth}
        {\small\textbf{\faBrain\ Extended Thinking}}
        \vspace{0.1cm}

        \terminal{0.95\columnwidth}{claude-code-py}{%
            \color{promptgreen}\textbf{You:} \color{white}think deep: wzorce\\
            projektowe w~tym projekcie?\\
            \separator\\[2pt]
            \infoicon\color{white}Thinking level: DEEP (20{,}000)\\[2pt]
            \color{codegray}$\therefore$ Thinking complete (2s, 120)\\[1pt]
            \toolicon Executing tool:\color{white}bash\\
            \toolok\color{white}Tool result (success)\\[1pt]
            \color{codegray}$\therefore$ Thinking complete (3s, 176)\\[1pt]
            \toolicon Executing tool:\color{white}read\_file\\
            \toolok\color{white}Tool result (success)\\[1pt]
            \color{codegray}\textit{...12 rund narz\k{e}dzi...}\\[2pt]
            \color{cyan}\textbf{Claude:} \color{white}Factory, Registry,\\
            Chain of Responsibility,\\
            Strategy, Pipeline, DI...%
        }

        \vspace{0.15cm}
        \footnotesize
        Dynamiczny limit rund narz\k{e}dzi:\\
        BASIC~5, STANDARD~8, DEEP~12, ULTRA~15.\\
        Graceful fallback przy limicie.

        \column{0.48\textwidth}
        {\small\textbf{\faUsersCog\ Wyspecjalizowani agenci}}
        \vspace{0.1cm}

        \terminal{0.95\columnwidth}{claude-code-py}{%
            \color{promptgreen}\textbf{You:} \color{white}Zr\'ob code review tego kodu\\
            \separator\\[2pt]
            \color{codegray}$\therefore$ Thinking...\\
            \color{codegray}$\therefore$ Thinking complete (3s, 420 tokens)\\[2pt]
            \toolicon Executing tool: \color{white}read\_file\\
            \toolok\color{white}Tool result (success)\\[2pt]
            \color{cyan}\textbf{Claude:} \color{white}Przegl\k{a}d kodu:\\
            \color{codered}\textbf{!} \color{white}SQL injection (linia 23)\\
            \color{codeyellow}\textbf{!} \color{white}N+1 query (linia 45)\\
            \toolok\color{white}Sp\'ojne nazewnictwo%
        }

        \vspace{0.15cm}
        \footnotesize
        Router wewnętrznie wybiera agenta\\
        (keyword matching + confidence).\\
        UI pozostaje przezroczyste.
    \end{columns}
\end{frame}

% ══════════════════════════════════════════════════════════════════════════
% SLAJD 6 — Architektura
% ══════════════════════════════════════════════════════════════════════════
\begin{frame}{Architektura}
    \begin{columns}[T]
        \column{0.48\textwidth}
        \begin{center}
        \begin{tikzpicture}[
            arr/.style={-{Stealth[length=5pt]}, thick, gray}
        ]
            \node[draw=claudepurple, fill=purple!8, minimum width=5.5cm, minimum height=0.7cm, rounded corners=3pt, font=\scriptsize]  (ui)   at (0, 4.0) {\textbf{UI} --- Textual TUI};
            \node[draw=claudepurple, fill=purple!14, minimum width=5.5cm, minimum height=0.7cm, rounded corners=3pt, font=\scriptsize] (app)  at (0, 3.0) {\textbf{Aplikacja} --- Orkiestracja};
            \node[draw=claudepurple, fill=purple!20, minimum width=5.5cm, minimum height=0.7cm, rounded corners=3pt, font=\scriptsize] (core) at (0, 2.0) {\textbf{Rdzeń} --- ChatPipeline + ToolLoop};
            \node[draw=claudepurple, fill=purple!26, minimum width=5.5cm, minimum height=0.7cm, rounded corners=3pt, font=\scriptsize] (cli)  at (0, 1.0) {\textbf{Klient} --- API Client + Factory};
            \node[draw=claudepurple, fill=purple!32, minimum width=5.5cm, minimum height=0.7cm, rounded corners=3pt, font=\scriptsize] (ext)  at (0, 0.0) {\textbf{Infra} --- Auth, Config, Security};

            \draw[arr] (ui)  -- (app);
            \draw[arr] (app) -- (core);
            \draw[arr] (core) -- (cli);
            \draw[arr] (cli) -- (ext);
        \end{tikzpicture}
        \end{center}

        \column{0.48\textwidth}
        \begin{center}
        {\small\textbf{Przepływ danych}}
        \vspace{0.2cm}

        \begin{tikzpicture}[
            box/.style={draw=claudepurple, fill=claudepurple!8, rounded corners, minimum height=0.6cm, minimum width=2.2cm, align=center, font=\tiny},
            arr/.style={-{Stealth[length=4pt]}, thick, claudepurple},
            node distance=0.3cm
        ]
            \node[box] (input)  {Zapytanie};
            \node[box, below=of input] (pipe) {ChatPipeline};
            \node[box, below=of pipe] (api) {Anthropic API};
            \node[box, below=of api] (stream) {Streaming};
            \node[box, below=of stream] (tool) {Tool Loop (max 5)};
            \node[box, below=of tool] (render) {Odpowiedź w~TUI};

            \draw[arr] (input) -- (pipe);
            \draw[arr] (pipe) -- (api);
            \draw[arr] (api) -- (stream);
            \draw[arr] (stream) -- (tool);
            \draw[arr] (tool) -- (render);
            \draw[arr, dashed] (tool.west) to[out=180,in=180] node[left, font=\tiny] {wynik} (api.west);
        \end{tikzpicture}
        \end{center}
    \end{columns}
\end{frame}

% ══════════════════════════════════════════════════════════════════════════
% SLAJD 7 — Podsumowanie
% ══════════════════════════════════════════════════════════════════════════
\begin{frame}{Podsumowanie}
    \begin{columns}[T]
        \column{0.48\textwidth}
        \begin{block}{Zrealizowane funkcje}
            \begin{itemize}\setlength{\itemsep}{1pt}
                \item Pełnoekranowy TUI ze streamingiem
                \item Integracja z~Anthropic API
                \item Tool calling --- 6~narzędzi
                \item Extended thinking --- 4~poziomy
                \item System 4~agentów
                \item Dual auth (OAuth + API key)
                \item Tryb interaktywny + CLI
            \end{itemize}
        \end{block}

        \column{0.48\textwidth}
        \begin{block}{Możliwe rozszerzenia}
            \begin{itemize}\setlength{\itemsep}{1pt}
                \item Obsługa wielu modeli (Gemini, GPT)
                \item Persistent historia konwersacji
                \item Wtyczki / własne narzędzia
                \item Integracja z~Git
                \item Web UI
            \end{itemize}
        \end{block}
    \end{columns}

    \vspace{0.4cm}
    \begin{center}
        \Large\textbf{Dziękuję za uwagę}

        \vspace{0.2cm}
    \end{center}
\end{frame}

\end{document}
